%%%%%%%%%%%%%%%%%%%%%%%%%%%%%%%%%%%%%%%%%%%%%%%%%%%%%%%
% SECTION 11: METHODS (Online Methods)
%%%%%%%%%%%%%%%%%%%%%%%%%%%%%%%%%%%%%%%%%%%%%%%%%%%%%%%
% 
% INPUTS NEEDED: 
%   - Technical implementation details
%   - Framework/library versions
%   - Training procedures summary
%
% FIGURES: None
% TABLES: None
% DEPENDENCIES: 
%   - Results sections (must align)
%   - Supplementary Methods (detailed version)
%
%%%%%%%%%%%%%%%%%%%%%%%%%%%%%%%%%%%%%%%%%%%%%%%%%%%%%%%

\section*{Methods}\label{sec:methods}

%%=============================================================%%
%% METHOD 1: Automated variant-evidence extraction
%%=============================================================%%

\subsection*{Automated variant--evidence extraction}

We built literature curation around the Reason--Action--Conclude--Think (ReACT) paradigm in DSPy. PDFs are parsed with PyPDF2 (with omniOCR fallback for scanned/image-heavy files). A GPT-4.1 Curation Agent executes explicit reasoning chains (DSPy ChainOfThought, Predict, ReAct for tool calls) against a formal DSPy signature that maps full-text input to CIViC-compatible evidence items. A Validation Engine enforces four layers: (i)~schema checks via Pydantic; (ii)~field checks for HGVS, PMID, and NCT patterns; (iii)~cross-entity integrity between evidence and variant records; and (iv)~clinical-logic consistency between study design, evidence level, and therapeutic relationships. A composite confidence score (evidence strength, completeness, and extraction fidelity) routes low-confidence items to human review. The asynchronous Python implementation preserves full reasoning traces and adheres to HGVS/HGNC/ClinicalTrials.gov standards, with FAIR and HL7 FHIR-genomics compatibility.

%%=============================================================%%
%% METHOD 2: Natural-language-to-SQL query system
%%=============================================================%%

\subsection*{Natural-language-to-SQL query system}

To remove the SQL barrier, we generated 21,000 natural-language question--SQL pairs across four complexity tiers using templated prompts populated with real CIViC entities. Each SQL target passed multi-stage validation (syntax, schema, executability). We fine-tuned Llama-3 3B with LoRA adapters (parameter-efficient training on Fireworks AI) to translate questions into execution-ready SQL over the curated OncoCITE-KB. Training details are provided in Supplementary Methods.

%%=============================================================%%
%% METHOD 3: Knowledge-enhanced clinical interface
%%=============================================================%%

\subsection*{Knowledge-enhanced clinical interface}

We added a five-component interface: a Knowledge Management System that canonicalizes gene aliases, variant notations, disease terms, and therapy names; a Multi-Step Reasoning Pipeline that reformulates failed or under-specified queries (constraint relaxation, synonym substitution, intent reinterpretation, and query decomposition); an Enhanced Database Manager for optimization and failure analysis; an Interactive Dialogue layer for collaborative refinement; and a Transparency Layer that logs all decisions. A Clinical Interpretation Engine converts query results into hierarchical, guideline-aligned summaries (context, key findings, clinical relevance, next steps). Phase-3 behavior was assessed in simulated clinical scenarios with standardized rubrics.

%%%%%%%%%%%%%%%%%%%%%%%%%%%%%%%%%%%%%%%%%%%%%%%%%%%%%%%
% METHODS BREAKDOWN:
%
% METHOD 1: AUTOMATED EXTRACTION
%
% Framework: DSPy + ReACT paradigm
% PDF parsing: PyPDF2 (omniOCR fallback)
% LLM: GPT-4.1 Curation Agent
%
% DSPy modules used:
%   - ChainOfThought
%   - Predict
%   - ReAct (for tool calls)
%
% Four-layer validation:
%   (i)   Schema checks (Pydantic)
%   (ii)  Field checks (HGVS, PMID, NCT patterns)
%   (iii) Cross-entity integrity
%   (iv)  Clinical-logic consistency
%
% Confidence scoring:
%   - Evidence strength
%   - Completeness
%   - Extraction fidelity
%   → Routes low-confidence to human review
%
% Standards compliance:
%   - HGVS
%   - HGNC
%   - ClinicalTrials.gov
%   - FAIR
%   - HL7 FHIR-genomics
%
%%%%%%%%%%%%%%%%%%%%%%%%%%%%%%%%%%%%%%%%%%%%%%%%%%%%%%%

%%%%%%%%%%%%%%%%%%%%%%%%%%%%%%%%%%%%%%%%%%%%%%%%%%%%%%%
% METHOD 2: NL-TO-SQL
%
% Training data:
%   - 21,000 Q-SQL pairs
%   - Four complexity tiers
%   - Templated prompts + real CIViC entities
%
% Validation stages:
%   - Syntax
%   - Schema
%   - Executability
%
% Model:
%   - Llama-3 3B
%   - LoRA adapters
%   - Fireworks AI platform
%
% → Detailed hyperparameters in Supplementary Methods
%
%%%%%%%%%%%%%%%%%%%%%%%%%%%%%%%%%%%%%%%%%%%%%%%%%%%%%%%

%%%%%%%%%%%%%%%%%%%%%%%%%%%%%%%%%%%%%%%%%%%%%%%%%%%%%%%
% METHOD 3: KNOWLEDGE-ENHANCED INTERFACE
%
% Five components:
%   1. Knowledge Management System
%      - Gene aliases
%      - Variant notations
%      - Disease terms
%      - Therapy names
%
%   2. Multi-Step Reasoning Pipeline
%      - Constraint relaxation
%      - Synonym substitution
%      - Intent reinterpretation
%      - Query decomposition
%
%   3. Enhanced Database Manager
%      - Query optimization
%      - Failure analysis
%
%   4. Interactive Dialogue Layer
%      - Collaborative refinement
%
%   5. Transparency Layer
%      - Decision logging
%
% Clinical Interpretation Engine outputs:
%   - Context summary
%   - Key findings
%   - Clinical relevance
%   - Next steps
%
% Evaluation:
%   - Simulated clinical scenarios
%   - Standardized rubrics
%
%%%%%%%%%%%%%%%%%%%%%%%%%%%%%%%%%%%%%%%%%%%%%%%%%%%%%%%

%%%%%%%%%%%%%%%%%%%%%%%%%%%%%%%%%%%%%%%%%%%%%%%%%%%%%%%
% CROSS-REFERENCE WITH SUPPLEMENTARY METHODS:
%
% Detailed content in Supplementary Methods includes:
%   □ DSPy module configurations
%   □ LoRA hyperparameters (rank=8, alpha=16, dropout=0.05)
%   □ Optimizer settings (AdamW, lr=2e-5, batch=16)
%   □ Training platform (A100 GPUs)
%   □ Dictionary sizes and sources
%   □ Trie data structure implementation
%   □ Fuzzy matching algorithms
%
%%%%%%%%%%%%%%%%%%%%%%%%%%%%%%%%%%%%%%%%%%%%%%%%%%%%%%%
